\section{Introduction}
\label{sec:introduction}
\paragraph{1.1}
The unobserved productive ceiling is not a neutral engineering constraint; it is the central macroeconomic variable linking capital accumulation, profitability, and class struggle. While actual output is driven by effective demand, productive capacity reflects the institutional limits of exploitation and the physical intensity of the labor process. \citet{Shaikh2016} offers a portable, accounting-based solution to measure this ceiling via the long-run relation between aggregate output and capital stock. However, treating the estimated output–capital coefficient as a stable, physical law of production risks masking the historical and distributive conflicts that govern the workplace. This chapter critically replicates and extends Shaikh’s strategy to determine whether this empirical coefficient represents a stable technical relation or a historically contingent reflection of class struggle.

\paragraph{1.2}
Distinct from the baseline approach of \citet{Shaikh2016}, this chapter reframes the aggregate output–capital coefficient as a capacity transformation elasticity ($\theta$). We define the latent capacity ceiling in log-linear form as:
\begin{equation}
\ln Y^p_t = \theta \ln K_t + \ln R^n_0,
\end{equation}
where $Y^p_t$ is potential capacity, $K_t$ is the capital stock, and $R^n_0$ is a scale constant. Under the standard balanced-growth assumption of post-Keynesian models, this elasticity is constrained to unity ($\theta = 1$), meaning capacity expands in lockstep with capital. If $\theta \neq 1$, however, capacity formation decouples from accumulation. An elasticity below unity ($\theta < 1$) characterizes a structural overaccumulation regime where capital growth outpaces capacity expansion, while $\theta > 1$ leads to chronic excess capacity. By analyzing $\theta$ directly, we move beyond Shaikh’s focus on short-run residual fluctuations to investigate the long-run structural stability of the output–capital relation.

\paragraph{1.3}
Evaluating this relationship reveals that Shaikh's baseline results are approximately recoverable, yet the underlying elasticity is highly sensitive to model selection. In replicating the single-equation ARDL model of \citet{Shaikh2016} using US corporate sector data (1947–2011), we come close to but do not exactly reproduce the original estimates due to undocumented lag structures and price deflator choices. Interrogating this fragility across a combinatoric specification space of 500 ARDL models reveals "disciplined non-uniqueness." Rather than converging to a single technical parameter, the estimated elasticity ranges from 0.65 to 0.95. This parameter volatility demonstrates that the recovered capacity path is a statistical artifact of researcher-chosen lag lengths and historical dummy variables, rather than a unique physical law.

\paragraph{1.4}
Shifting the analysis to a joint system framework reveals that the bivariate output–capital relationship is spurious. In a bivariate Vector Error Correction Model (VECM), the output–capital relation fails cointegration and dynamic stability tests across all lag profiles. Cointegration and system-level stability emerge only when the logged rate of exploitation ($e_t$) enters the state vector in a trivariate VECM. This cointegrating survival provides the central empirical finding of the chapter: the rate of exploitation is not a neutral accounting ratio, but a proxy for the balance of class power that dictates shop-floor discipline, machine speed-ups, and the choice of technique. The empirical requirement of $e_t$ in the cointegrating vector indicates that capacity utilization is a distributionally conditioned phenomenon, requiring explicit class struggle controls to establish a stable long-run attractor.

\paragraph{1.5}
These findings define the historical limits of the accounting-based capacity measurement of \citet{Shaikh2016}. While the output–capital ratio remains a useful anchor to construct capacity utilization, the resulting ceiling cannot be treated as an insulated engineering benchmark. Cointegration survives only when the model is conditioned by the rate of exploitation and stabilized by historical dummy controls that mark structural breaks in post-war US accumulation. Ultimately, the long-run output–capital relation cannot be estimated in isolation from the institutional conflicts of the capitalist labor process, which physically govern how capital is converted into potential output.

\paragraph{1.6}
To establish this argument, the chapter unfolds along a narrative that historicizes the measurement problem before executing the econometric tests. We first reconstruct the historical evolution of capacity metrics in Section 2, tracing how survey indicators and output-gap routines shifted from tracking industrial output to enforcing price-level and wage discipline. In Section 3, we develop the conceptual Leontief framework, showing how omitting distribution introduces an identity trap that biases the capital coefficient, and detailing how \citet{Basu2022}'s viability condition governs endogenous technical change. Section 4 presents the staged empirical stress tests—reconstructing the baseline, mapping the ARDL specification space, and testing system-level VECM cointegration. Section 5 concludes by detailing the theoretical implications of structural overaccumulation ($\hat{\theta} < 1.0$) for profit rate crisis theory, evaluating the interpretation of these distributive coefficients, and outlining the Chapter 2 point of departure for Global South economies.
